\documentclass[11pt]{article}
\usepackage{arxiv-preprint}

% =========================================================================
% DRAFT MARKERS
%   \todo{...}   : missing content or numbers pending an experiment run
%   \claim{...}  : an empirical claim that a reader should be able to
%                  trace back to a JSON in publish/01-fenrir/code/outputs/
% =========================================================================
\newcommand{\todo}[1]{{\bfseries\color{red}[TODO: #1]}}
\newcommand{\claim}[1]{#1}
% escape helper: use \us for underscore inside \todo to avoid math-mode parse
\newcommand{\us}{\textunderscore}

\title{FENRIR: A Fast-Weight Mixer for Single-Pass Chain Composition}

\author[1]{Voidwolf}
\affil[1]{Independent researcher \\ \texttt{research@voidwolf.space}}

\date{\today}

\begin{document}
\maketitle

\begin{abstract}
We study when and how a fast-weight mixer with a state-dependent write
address learns to compose 2-hop retrieval chains inside a single
forward pass. The mixer, which we call FENRIR, computes an eager
lookup $a_t = M_{t-1} k_t$ into the current state, adds it to the
current key to form $k_{\text{eff}} = k_t + a_t$, and writes the outer
product $\beta_t k_{\text{eff}} v_t^{\top}$ into the fast-weight state.
On a synthetic 2-hop chain-composition task at $K_1 \in \{1, \dots, 4\}$
with $K_2 = 4$ candidate values, we train up to $n = 10$ seeds per
variant under matched protocol and report four mechanism-level
findings: (a) chain composition to $K_1 = 4$ inside a single forward
pass is attainable, but formation is bimodal: under mixed-mode
training \claim{$5/10$} \texttt{rev} seeds cross the ladder to
$\text{acc} \geq 0.895$ while the other $5/10$ stall at chance;
under joint-mode training \claim{$6/10$} pass at $\text{acc} \geq
0.958$; the \texttt{fwd} operator direction has a much lower per-seed
pass probability ($0/5$ mixed; $1/5$ joint at $0.977$) but is not
categorically failing under joint mode;
(b) on passing seeds the answer becomes decodable from the third
layer's residual at accuracy matching training, while failing seeds
show no localisation at any layer; (c) an inference-time knockout of
the eager address-perturbation on trained checkpoints collapses
passing-seed accuracy from $\claim{\sim 0.96}$ to $\claim{\sim 0.03}$
(worse than chance), with \texttt{fwd} checkpoints collapsing to
$0.000$, showing the term is not incidental but the computational
pathway the model routes through; (d) a mid-scan probe traces this
pathway to a buildup inside the state during the hop-2 write window,
and zeroing the eager term at probe time collapses the buildup
consistent with (c). We compare two operator directions (reverse-lookup
$k + M k$ vs forward-lookup $k + M^{\top} k$) and report their
empirical envelopes side by side. We make no strong claim of
primitive-level novelty; the paper's contribution is the
mechanistic characterisation and the audit-safe reproduction
workflow. All code, tests, and mechanism probes are packaged as a
standalone repository, and every number in the paper is traceable to
a committed JSON.
\end{abstract}

\section{Introduction}
\label{sec:intro}

A \emph{chain composition} step follows two facts through their shared
intermediate: given $(a, b)$ and $(b, c)$ known, retrieve $c$ from $a$.
Standard fast-weight and delta-family mixers store each fact
independently and depend on stacking multiple layers, or multiple
forward passes, to compose across them. The mechanism question we ask
in this paper is narrower: given a mixer whose write address itself
depends on the current state, when and how does it learn to compose a
2-hop chain inside a single forward pass?

The primitive we study, which we call FENRIR, is an address-perturbed
fast-weight mixer with two operator directions we treat as sibling
variants: \emph{reverse-lookup} (\texttt{rev}) with
$k_{\text{eff}} = k + M k$ and \emph{forward-lookup} (\texttt{fwd})
with $k_{\text{eff}} = k + M^{\top} k$. Both share the same
outer-product write and the same forward read; both belong to the
outer-product / delta-family lineage (\S\ref{sec:related}). We do not
claim the primitive itself is new. Prior fast-weight work has explored
many state-dependent write variants and it is plausible that the
FENRIR update rule appears somewhere under a different name; we hedge
this in \S\ref{sec:related} and give the closest prior art we found.

\textbf{What the paper does claim} is a mechanism-level account of when
composition forms in this primitive and what it looks like inside the
model when it does. Concretely:

\begin{enumerate}
\item \textbf{Bimodal formation (\S\ref{sec:results:bimodality}).}
      Under a matched-protocol $n=5$-seed run at $K_1 = 4$,
      under mixed-mode training \claim{$5/10$} \texttt{rev} seeds
      cross the full K1 ladder to $\text{acc} \geq 0.895$; under
      joint-mode training \claim{$6/10$} pass at $\text{acc} \geq
      0.958$. The failing cluster in both curricula lands at chance;
      passing and failing seeds sit at cleanly separated clusters.
      The \texttt{fwd} operator direction has a much lower per-seed
      pass probability ($0/5$ mixed; $1/5$ joint at $0.977$), showing
      curriculum can partially rescue \texttt{fwd} formation. A
      single-seed report would report one number from that
      distribution and be misleading.

\item \textbf{Layer localisation (\S\ref{sec:results:mechanism}).} On
      passing seeds a logit-lens probe shows the answer becoming
      decodable at the third layer's residual output at accuracy
      matching training. On failing seeds no layer localises anything.
      The distinction between passing and failing is not distributed;
      it is a specific layer's behaviour.

\item \textbf{Inference-time causal test
      (\S\ref{sec:ablations:inference-ko}).} On trained checkpoints,
      forcing the mixer's eager term to zero at every position of the
      forward pass collapses passing-\texttt{rev} accuracy from
      $\text{acc} \sim 0.96$ to $\text{acc} \sim 0.03$ (worse than
      chance), and collapses \texttt{fwd} checkpoints to $0.000$ hits.
      The eager address-perturbation is not incidental; it is the
      computational pathway the model routes through.

\item \textbf{Mid-scan mechanism (\S\ref{sec:results:midscan}).} A
      per-position probe of the mixer state at the layer where
      composition is constructed shows the answer signal building up
      inside the state during the hop-2 write window; zeroing the
      eager term at probe time collapses the buildup, consistent with
      the inference-time result above.

\item \textbf{Audit workflow (\S\ref{sec:setup:audit}).} The mechanism
      probes read from the mixer's own cache or use the mixer's own
      eager operator method, so an external reconstruction cannot
      silently compute a different operator than the mixer used. Every
      convention claim (M storage direction, eager operator per
      variant, read direction) is expressed as a pytest test; the
      repository is a working example of an audit workflow that
      catches the operator-convention bug class described in
      Appendix~\ref{app:audit}.
\end{enumerate}

The paper's tables read from JSON files committed under
\texttt{code/outputs/}. Every number in the text is traceable to one.
Given the formation-mode bimodality above, the reader should treat all
per-seed numbers as the raw ingredient, and any per-variant summary
statistic as an aggregate over that distribution rather than a
representative point estimate.

\section{Background}
\label{sec:background}

\textbf{Fast-weight programmers.} A linear-attention layer can be
rewritten as a fast-weight programmer whose associative memory
$M \in \mathbb{R}^{d \times d}$ is updated by outer-product writes and
read by matrix-vector products \citep{schmidhuber1992fastweights,
schlag2021linear}. Each per-position write $M \mathrel{{+}{=}} k_t
v_t^{\top}$ stores an association between a projected key and value;
each read $y_t = M^{\top} q_t$ retrieves values by matching the query
against stored keys.

\textbf{Delta-family updates.} A pure additive write accumulates
without bound and interferes as the state fills. The delta rule fixes
this by writing a residual against the current prediction:
$M \mathrel{{+}{=}} \beta_t \, k_t \, (v_t - M^{\top} k_t)^{\top}$
\citep{schlag2021deltanet}. \citet{yang2024parallelizing} showed the
delta rule admits a chunk-parallel form via a triangular closure,
making delta-family mixers tractable at LM scale;
\citet{yang2024gated} added scalar gating.

\textbf{Chain composition and associative recall.} Small-scale
compositional-recall tasks are a common substrate for probing mixer
capabilities. Zoology and BASED
\citep{arora2023zoology,arora2024based} study single-hop associative
recall in linear attention and gated recurrent architectures.
Grokked-Transformers \citep{wang2024grokked} shows dense
transformers acquire multi-hop compositional recall only under
extended training, framing chain composition as a capability gap
worth closing with better architectural inductive biases.

\textbf{Mechanism probes.} Logit-lens
\citep{nostalgebraist2020logitlens} reads intermediate residuals
through the model's own LM head to locate where a target token
becomes decodable. Induction-head style circuit work
\citep{olsson2022incontext} makes causal claims about intermediate
representations via targeted knockouts. Our mid-scan probe extends
this pattern to per-timestep reads of the fast-weight state.

\textbf{State-space scaffolds.} The FENRIR block wraps its mixer in a
Mamba \citep{gu2023mamba,dao2024mamba2} S4-style scaffold
(\texttt{in\_proj}, causal depthwise conv, $z$-gate, output
projection), so the comparison against Mamba as a non-fast-weight
baseline is a swap-in change to the mixer only.

\section{Method}
\label{sec:method}

\subsection{The FENRIR mixer}
\label{sec:method:mixer}

At each position $t$, given projections $q_t, k_t, v_t \in \mathbb{R}^d$
and a scalar gate $\beta_t \in (0, 1)$, the mixer maintains a
fast-weight state $M \in \mathbb{R}^{d \times d}$ under the storage
convention $M[i, j] = k_{\text{eff}}[i] \cdot v[j]$ (keys index axis
$0$, values index axis $1$). The per-position update is
\begin{align}
a_t &= \begin{cases}
        M_{t-1}\, k_t             & \text{(rev variant)}\\
        M_{t-1}^{\top}\, k_t      & \text{(fwd variant)}
       \end{cases}\\
k_{\text{eff},t} &= k_t + a_t \\
M_t &= M_{t-1} + \beta_t \, k_{\text{eff},t} v_t^{\top} \\
y_t &= M_t^{\top} q_t
\end{align}

The two variants differ only in the direction of the eager
address-perturbation $a_t$: \texttt{rev} matches $k_t$ against stored
values and returns stored keys; \texttt{fwd} matches $k_t$ against
stored keys and returns stored values. The write and the read are
identical across variants.

Both variants require $d_{\text{key}} = d_{\text{value}}$ so
$k_{\text{eff},t}$ type-checks. The mixer's forward pass uses a
sequential per-position loop. For \texttt{rev} we additionally provide a
chunk-parallel kernel that reduces each chunk to a single unit-lower-
triangular solve; the derivation is standard and appears in the
appendix.

\subsection{Scaffold}

Each FENRIR block wraps the mixer in a Mamba-style S4 scaffold:
\texttt{in\_proj} $\to$ depthwise causal conv1d $\to$ SiLU $\to$
$z$-gate $\to$ mixer $\to$ readout $\to$ elementwise gate by
$\mathrm{SiLU}(z)$ $\to$ \texttt{out\_proj}. Blocks are stacked with
pre-norm RMS residuals; the final residual reads out through a tied LM
head.

\section{Experimental setup}
\label{sec:setup}

\subsection{Task}
\label{sec:setup:task}

We use a 2-hop injective chain-composition task with a
\emph{truncated}-answer supervision protocol. Each episode has:
\begin{itemize}
\item $K_1$ hop-1 facts of the form $(e_i, b_i)$ (entity $\to$ bridge);
\item $K_2$ hop-2 facts of the form $(c_j, v_j)$ (bridge-key $\to$
      value); the bridges $\{b_i\}$ are an \emph{injective} subset of
      $\{c_j\}$ enforced by \texttt{randperm};
\item a query $q = e_{\text{target}}$;
\item the answer $v_{f(\text{target})}$ where $f: [K_1] \to [K_2]$ is
      the injective map $b_i = c_{f(i)}$.
\end{itemize}

The sequence layout is
\texttt{FACT $e_1 b_1$ SEP $e_2 b_2$ SEP \dots\ SEP $c_1 v_1$ SEP
$c_2 v_2$ \dots\ QTOK $q$ ATOK} and the model is teacher-forced up to
and including \texttt{ATOK}; loss is cross-entropy against the answer
token at the \texttt{ATOK} position. Chance is $1 / K_2 = 0.25$ for
$K_2 = 4$. The injective constraint blocks a ``collision luck'' shortcut
that inflates chain accuracy on the non-injective version of this task.

\subsection{Model configuration}

Unless noted, $d_{\text{model}} = 256$, $d_{\text{key}} = d_{\text{value}} = 32$,
$3$ FENRIR blocks, \texttt{expand} $= 2$, conv kernel $= 4$. Vocabulary is
\claim{$263$ tokens} ($7$ control tokens $+$ $128$ entities $+$ $128$
values). Parameter count is \claim{$1{,}453{,}827$}.

\subsection{Training protocol}
\label{sec:setup:training}

Optimizer is AdamW with $\text{lr} = 3\mathrm{e}{-4}$, $\text{betas} =
(0.9, 0.95)$, $\text{weight\_decay} = 0.01$. LR schedule is cosine with
$200$-step linear warmup and floor $3\mathrm{e}{-5}$. Precision is bf16
autocast on CUDA. Gradient clip is $1.0$ (global norm). Batch size
$= 64$. Two training curricula are compared:
\begin{itemize}
\item \textbf{mixed} (Table~\ref{tab:ladder-mixed}): at each step, cycle
      through $K_1 \in \{1, \dots, 4\}$. $4000$ steps total.
\item \textbf{joint} (Table~\ref{tab:ladder-joint}): at each step,
      sample one batch each at $K_1 = 1$ and $K_1 = 4$, average the
      cross-entropy losses. $5000$ steps total.
\end{itemize}

Optimizer hyperparameters follow the large-scale LM convention (GPT-3,
Chinchilla, LLaMA) rather than PyTorch's Adam defaults ($\beta_2 =
0.95$ rather than $\beta_2 = 0.999$). At this experiment's scale
($\sim 1.5$M parameters, $\lesssim 5000$ steps) the choice has limited
effect. We did not ablate it.

\subsection{Reference hardware}

All numbers reported here were produced on an NVIDIA RTX~5090 (32~GB,
CUDA 13.3, PyTorch~2.x, pinned as \texttt{torch>=2.1} in the reproduction
repository's \texttt{requirements.txt}) under bf16 autocast.

\subsection{Audit as code}
\label{sec:setup:audit}

Every claim about the mixer's storage convention or its eager operator
per variant is expressed as a pytest test in
\texttt{code/tests/test\_conventions.py}. Convention drift in the code
would break CI. In particular, the mixer stores its per-step eager term
under a variant-agnostic key so a downstream probe cannot silently
compute a different operator than the mixer used (an earlier revision
of this line of work had exactly this failure mode).

A publisher-side parity script (not shipped with the reproduction
repository) confirms the fresh implementation is byte-identical to the
project's rig-of-record reference implementations for both variants,
and byte-identical to the reference task sampler across the $K_1$/$K_2$
configurations we train on.

\subsection{Mechanism probes}
\label{sec:setup:probes}

We use three probes, all variant-aware and cache-safe:

\begin{itemize}
\item \textbf{F3 (logit lens).} At each layer's residual output, apply
      the model's own \texttt{RMSNorm} and LM head; report the fraction
      of episodes whose \texttt{ATOK} residual decodes the answer token
      as the argmax restricted to the $K_2$ candidate value tokens.
\item \textbf{F4 (mid-scan).} At a chosen layer, cache
      $(q, k, v, \beta)$ during the forward pass. Externally
      reconstruct the state $M_t$ at every timestep using the mixer's
      \emph{own} eager operator method (so the probe cannot diverge from
      the mixer's semantic). Compute
      $y'_t = q_{\text{ATOK}}^{\top} M_t$ at every $t$ and fit a
      $K_2$-way constrained logistic regression to
      $y'_t \to \{\text{bridge},\text{answer}\}$. An optional knockout
      arm zeros the eager term during the reconstruction.
\item \textbf{F5 ($Q$/$K$ alignment).} Three cosine measurements each
      with a same/control setup: (a) $\cos(q_{\text{ATOK}},
      k[k_1^{\text{target}}])$, (b) $\cos(k[b_i], k[c_{f(i)}])$,
      (c) $\cos(v[b_i], k[e_i])$. Reports the mean over $n$ batches.
\end{itemize}

\section{Results}
\label{sec:results}

\subsection{K-hop ladder}
\label{sec:results:ladder}

\begin{table}[h]
\centering
\caption{Chain-composition accuracy on the K1 ladder under \emph{mixed}
training (4000 steps, $n = 10$ seeds for \texttt{rev}, $n = 5$ for
\texttt{fwd}). Chance $= 0.25$. Numbers are per-seed final accuracies
at each rung; passing seeds (final $K_1{=}4$ acc $\geq 0.85$) are
bolded. The one borderline seed (\texttt{rev} s8 at $0.895$) sits
closer to the passing cluster than the failing one.}
\label{tab:ladder-mixed}
\begin{tabular}{llcccc}
\toprule
Variant & Seed & $K_1 = 1$ & $K_1 = 2$ & $K_1 = 3$ & $K_1 = 4$ \\
\midrule
\texttt{rev} & 0 & 0.996 & 0.992 & 0.965 & \textbf{0.929} \\
\texttt{rev} & 1 & 0.994 & 0.977 & 0.955 & \textbf{0.914} \\
\texttt{rev} & 2 & 0.902 & 0.420 & 0.311 & 0.260 \\
\texttt{rev} & 3 & 0.970 & 0.503 & 0.335 & 0.257 \\
\texttt{rev} & 4 & 0.877 & 0.427 & 0.314 & 0.255 \\
\texttt{rev} & 5 & 0.898 & 0.460 & 0.310 & 0.262 \\
\texttt{rev} & 6 & 0.977 & 0.481 & 0.312 & 0.273 \\
\texttt{rev} & 7 & 0.995 & 0.979 & 0.962 & \textbf{0.942} \\
\texttt{rev} & 8 & 0.991 & 0.970 & 0.939 & \textbf{0.895} \\
\texttt{rev} & 9 & 0.992 & 0.975 & 0.944 & \textbf{0.917} \\
\midrule
\texttt{fwd} & 0 & 0.996 & 0.476 & 0.329 & 0.265 \\
\texttt{fwd} & 1 & 0.761 & 0.508 & 0.376 & 0.328 \\
\texttt{fwd} & 2 & 0.819 & 0.401 & 0.305 & 0.253 \\
\texttt{fwd} & 3 & 0.998 & 0.493 & 0.330 & 0.273 \\
\texttt{fwd} & 4 & 0.991 & 0.492 & 0.342 & 0.263 \\
\bottomrule
\end{tabular}
\end{table}

Under mixed-mode training, the \texttt{rev} variant passes the full
ladder to $K_1{=}4$ on \claim{$5/10$} seeds at $\text{acc} \geq 0.895$
(with the borderline s8 at $0.895$ closer to the passing cluster than
the failing one) and stalls at chance on the other $5/10$. The
\texttt{fwd} variant does not pass on any seed under this curriculum,
plateauing at \claim{$\text{acc} \lesssim 0.33$} at $K_1 = 4$; we did
not extend \texttt{fwd} to $n{=}10$ because the $n{=}5$ signal is
already unambiguous.

\begin{table}[h]
\centering
\caption{Same as Table~\ref{tab:ladder-mixed} but under \emph{joint}
training (5000 steps, sampling one batch each at $K_1{=}1$ and $K_1{=}4$
per step, averaged CE). Joint mode dramatically improves the \texttt{rev}
pass rate; \texttt{fwd} remains at chance.}
\label{tab:ladder-joint}
\begin{tabular}{llcccc}
\toprule
Variant & Seed & $K_1 = 1$ & $K_1 = 2$ & $K_1 = 3$ & $K_1 = 4$ \\
\midrule
\texttt{rev} & 0 & 0.999 & 0.963 & 0.975 & \textbf{0.971} \\
\texttt{rev} & 1 & 0.999 & 0.977 & 0.979 & \textbf{0.983} \\
\texttt{rev} & 2 & 1.000 & 0.463 & 0.271 & 0.254 \\
\texttt{rev} & 3 & 0.999 & 0.988 & 0.988 & \textbf{0.984} \\
\texttt{rev} & 4 & 0.999 & 0.961 & 0.971 & \textbf{0.958} \\
\texttt{rev} & 5 & 0.998 & 0.418 & 0.287 & 0.250 \\
\texttt{rev} & 6 & 1.000 & 0.938 & 0.973 & \textbf{0.966} \\
\texttt{rev} & 7 & 1.000 & 0.946 & 0.988 & \textbf{0.979} \\
\texttt{rev} & 8 & 0.999 & 0.408 & 0.278 & 0.265 \\
\texttt{rev} & 9 & 0.999 & 0.421 & 0.290 & 0.265 \\
\midrule
\texttt{fwd} & 0 & 0.999 & 0.457 & 0.280 & 0.250 \\
\texttt{fwd} & 1 & 0.999 & 0.453 & 0.278 & 0.262 \\
\texttt{fwd} & 2 & 0.999 & 0.413 & 0.268 & 0.256 \\
\texttt{fwd} & 3 & 0.999 & 0.975 & 0.978 & \textbf{0.977} \\
\texttt{fwd} & 4 & 0.999 & 0.408 & 0.279 & 0.249 \\
\bottomrule
\end{tabular}
\end{table}

Under joint-mode training, the \texttt{rev} variant passes the full
ladder on \claim{$6/10$} seeds at $\text{acc} \geq 0.958$, versus
$5/10$ (with one borderline at $0.895$) under mixed-mode. Under both
curricula the failing cluster lands cleanly at chance ($\text{acc}
\leq 0.27$), preserving a sharp bimodal formation pattern.

\textbf{Curriculum rescues \texttt{fwd} formation partially.} A
notable finding under joint-mode is that the \texttt{fwd} variant,
which showed $0/5$ pass at $K_1=4$ under mixed-mode, has \claim{$1/5$}
pass under joint at $\text{acc} = 0.977$ (seed~3). This falsifies the
reading that \texttt{fwd} fails categorically; the operator-direction
asymmetry is better characterised as \texttt{fwd} having a much lower
per-seed pass probability under matched protocol than \texttt{rev}
($1/5$ vs. $6/10$ under joint), one that the joint curriculum
partially rescues. The passing \texttt{fwd} trajectory crystallises
later than passing \texttt{rev} seeds (rung $K_1{=}4$ reaches
$\geq 0.9$ near step 3200 for \texttt{fwd}~s3 vs. steps 3000--4200
for the passing \texttt{rev} seeds).

\subsection{Formation-mode bimodality}
\label{sec:results:bimodality}

The pass-vs-fail split in Table~\ref{tab:ladder-mixed} is not seed noise
in the ordinary sense: passing seeds land at $\text{acc} \geq 0.91$ at
$K_1 = 4$, and failing seeds at $\text{acc} \leq 0.27$. The $K_1{=}2$
column also splits sharply between two clusters ($\geq 0.98$ vs
$\lesssim 0.51$). This bimodal ``formation'' pattern is consistent
with what earlier work on this mixer family reported anecdotally at
$n{=}1$--$3$; the $n{=}5$ sample here quantifies the pass rate directly.

The training-time signature separates the two modes well before
convergence. The first ladder rung ($k_1$) saturates near $1.0$ in
essentially every run, passing and failing alike, so it carries no
predictive signal on its own. The distinguishing event is the lift of the
second rung: in passing seeds $k_2$ begins climbing above the $1/K_2$ floor
partway through training and drags $k_3$ and $k_4$ up behind it, whereas
failing seeds hold $k_2$ through $k_4$ at chance for the entire run. Under
joint-mode training the passing seeds cross $\text{acc} = 0.95$ between
roughly step $3000$ and step $4400$ (recorded per run as the
\texttt{step\_to\_95} field in the committed result JSONs), while failing
seeds never reach it. The fate of a run is thus legible from the $k_2$
trajectory long before the final step, but not from $k_1$.

\subsection{Mechanism localisation (F3)}
\label{sec:results:mechanism}

\begin{table}[h]
\centering
\caption{F3 logit-lens ATOK-decodable accuracy at each layer's residual,
chance $= 0.25$. Passing seeds show the answer becoming decodable at
one of the final two residuals; failing seeds show no localisation
anywhere. Bolded columns mark where each passing seed first exceeds
$0.85$. Note the mechanism-class split within passing seeds:
some compose at $L_1$ (final layer), some already at $L_2$
(penultimate).}
\label{tab:f3}
\begin{tabular}{llcccc}
\toprule
Variant & Seed (curriculum) & \texttt{embed} & after $L_0$ & after $L_1$ & after $L_2$ \\
\midrule
\multicolumn{6}{l}{\emph{passing seeds -- compose at final layer $L_2$}} \\
\texttt{rev} & 0 (mixed) & 0.250 & 0.264 & 0.358 & \textbf{0.944} \\
\texttt{rev} & 1 (mixed) & 0.271 & 0.259 & 0.259 & \textbf{0.926} \\
\texttt{rev} & 0 (joint) & 0.249 & 0.323 & 0.326 & \textbf{0.973} \\
\texttt{rev} & 1 (joint) & 0.271 & 0.275 & 0.257 & \textbf{0.988} \\
\texttt{rev} & 3 (joint) & 0.256 & 0.258 & 0.254 & \textbf{0.990} \\
\texttt{rev} & 4 (joint) & 0.238 & 0.237 & 0.271 & \textbf{0.962} \\
\midrule
\multicolumn{6}{l}{\emph{passing seeds -- compose at penultimate layer $L_1$}} \\
\texttt{rev} & 7 (mixed) & 0.249 & 0.257 & \textbf{0.931} & 0.932 \\
\texttt{rev} & 8 (mixed) & 0.250 & 0.453 & \textbf{0.898} & 0.902 \\
\texttt{rev} & 9 (mixed) & 0.261 & 0.513 & \textbf{0.936} & 0.936 \\
\texttt{rev} & 6 (joint) & 0.255 & 0.261 & \textbf{0.968} & 0.975 \\
\texttt{rev} & 7 (joint) & 0.249 & 0.256 & \textbf{0.983} & 0.983 \\
\texttt{fwd} & 3 (joint) & 0.256 & 0.249 & \textbf{0.981} & 0.981 \\
\midrule
\multicolumn{6}{l}{\emph{failing seeds -- no localisation}} \\
\texttt{rev} & 2 (joint) & 0.260 & 0.253 & 0.252 & 0.246 \\
\texttt{rev} & 5 (joint) & 0.257 & 0.246 & 0.255 & 0.252 \\
\texttt{fwd} & 0 (mixed) & 0.252 & 0.250 & 0.249 & 0.262 \\
\texttt{fwd} & 0 (joint) & 0.250 & 0.252 & 0.250 & 0.248 \\
\multicolumn{6}{c}{\ldots\ (rest of failing seeds omitted for brevity, see \texttt{code/outputs/table2\us{}\us{}*f3.json})} \\
\bottomrule
\end{tabular}
\end{table}

A finding not present in the original single-seed reports of this
mixer family is that passing seeds separate into two mechanism
classes: half compose at the final layer ($L_2$ decodes strongly,
$L_1$ does not) and half compose one layer earlier ($L_1$ already
decodes, $L_2$ matches). Both classes achieve comparable
$K_1{=}4$ accuracy. Whether the mixer forms its composition at the
penultimate or final layer appears to be a seed-dependent choice
uncorrelated with pass/fail. The one passing \texttt{fwd} seed also
lands in the penultimate-layer class.

\subsection{Mid-scan buildup (F4)}
\label{sec:results:midscan}

At $L_1$ on the passing \texttt{rev}~s0 checkpoint, the answer signal
$q_{\text{ATOK}}^{\top} M_t$ traces the following trajectory across
positions (chance $= 0.25$):

\begin{center}
\begin{tabular}{lcccccccccc}
\toprule
position & \dots & $k_2[0]$ & $v_0$ & $k_2[1]$ & $v_1$ & $k_2[2]$ & $v_2$ & $k_2[3]$ & $v_3$ & \texttt{ATOK} \\
\midrule
answer   & \dots & 0.27     & 0.27  & 0.33     & 0.32  & 0.47     & 0.47  & 0.56     & 0.64  & \textbf{0.82} \\
\bottomrule
\end{tabular}
\end{center}

The answer is at chance throughout the hop-1 write window, remains near
chance early in the hop-2 window, and climbs monotonically through the
last three value writes to $0.82$ at \texttt{ATOK}. Bridge signal shows
a symmetric pattern in the hop-1 window (details in the appendix).

\textbf{Causal test.} Zeroing the eager term $a_t = M_{t-1} k_t$ during
the reconstruction (F4 with \texttt{--eager-knockout}) collapses the
\texttt{ATOK} answer signal on the same checkpoint to
\claim{$0.267$ (chance)}. The address-perturbation is causally
necessary for the mid-scan buildup.

\textbf{F4 sweep across representative checkpoints.} Table
\ref{tab:f4-sweep} extends the F4 probe to 10 checkpoints at $L_1$ and
adds an $L_0$ probe on the three passing seeds whose F3 decode only
appeared at the final residual. Findings:

\begin{itemize}
\item \textbf{$L_1$ composers} (4 passing seeds): F4 shows large
      buildup ($\sim$$0.94$) at $L_1$ and full collapse to chance
      under KO. Matches the F3 pattern where the answer decodes at
      $L_2$ (the residual immediately after $L_1$'s mixer).
\item \textbf{$L_0$ composer} (rev joint s0): F4 buildup shifts one
      block earlier -- $L_0$ base $= 0.965$, $L_0$ KO $= 0.280$ --
      and $L_1$ shows no buildup. F3 decodes only at the final
      residual, consistent with the answer being computed early but
      requiring two more residuals to reach the LM head.
\item \textbf{Diffuse composers} (rev joint s3, s4): Two passing
      seeds show no single-layer buildup at either $L_0$ or $L_1$
      (both at chance), yet decode strongly at the final F3 layer
      and achieve $K_1{=}4 \geq 0.98$. The mechanism is presumably
      spread across $L_2$ or distributed across layers; we did not
      resolve which because the $L_2$ F4 probe exhausted its LR
      convergence budget on multiple attempts.
\item \textbf{Failing seeds}: all show chance-level $L_1$ ATOK
      signal both with and without KO.
\end{itemize}

\begin{table}[h]
\centering
\caption{F4 mid-scan ATOK answer signal, baseline (eager on) vs.
\emph{eager-KO} (eager term zeroed during reconstruction). Chance
$= 0.25$. Passing seeds fall into three mechanism classes:
\emph{$L_1$ composers} show large buildup at $L_1$;
\emph{$L_0$ composers} show it at $L_0$ instead (with $L_1$ silent);
\emph{diffuse composers} show no single-layer buildup at either.
KO collapses the buildup to chance in the classes where it exists.}
\label{tab:f4-sweep}
\begin{tabular}{lllccc}
\toprule
Class & Variant/Seed & Probed & base & KO & Curriculum \\
\midrule
\multicolumn{6}{l}{\emph{$L_1$ composers -- F4 buildup at $L_1$; F3 decodes at $L_2$}} \\
                & \texttt{rev} s0 & $L_1$ & 0.917 & 0.267 & mixed \\
                & \texttt{rev} s7 & $L_1$ & 0.935 & 0.280 & mixed \\
                & \texttt{rev} s6 & $L_1$ & 0.958 & 0.262 & joint \\
                & \texttt{fwd} s3 & $L_1$ & 0.986 & 0.277 & joint \\
\midrule
\multicolumn{6}{l}{\emph{$L_0$ composer -- F4 buildup at $L_0$; F3 decodes only at final residual}} \\
                & \texttt{rev} s0 & $L_0$ & 0.965 & 0.280 & joint \\
                & \texttt{rev} s0 & $L_1$ & 0.296 & 0.309 & joint \\
\midrule
\multicolumn{6}{l}{\emph{Diffuse composers -- no single-layer buildup at $L_0$ or $L_1$; still pass}} \\
                & \texttt{rev} s3 & $L_0$ & 0.268 & 0.267 & joint \\
                & \texttt{rev} s3 & $L_1$ & 0.281 & 0.273 & joint \\
                & \texttt{rev} s4 & $L_0$ & 0.268 & 0.255 & joint \\
\midrule
\multicolumn{6}{l}{\emph{Failing seeds -- chance at $L_1$ base and KO}} \\
                & \texttt{rev} s2 & $L_1$ & 0.258 & 0.273 & joint \\
                & \texttt{rev} s5 & $L_1$ & 0.265 & 0.258 & joint \\
                & \texttt{fwd} s0 & $L_1$ & 0.277 & 0.254 & mixed \\
                & \texttt{fwd} s0 & $L_1$ & 0.262 & 0.256 & joint \\
\bottomrule
\end{tabular}
\end{table}

\subsection{Q/K alignment (F5)}
\label{sec:results:alignment}

We ran the $Q$/$K$ alignment probe on a small number of checkpoints
rather than the full sweep, so we report it as preliminary and
single-seed. On a passing \texttt{rev} seed (s0) at $L_2$,
prediction~(b), $\cos(k[b_i], k[c_{f(i)}])$, shows a positive gap over
its control ($0.474$ versus $0.436$, a $+0.038$ margin); on a failing
\texttt{rev} seed (s2) the same gap shrinks to $+0.012$. The gap is
larger at $L_1$, where the composition is being constructed rather than
read out. These point measurements are directionally consistent with the
F3 and F4 results (the aligned representation is present in passing seeds
and absent in failing ones), but we draw no quantitative conclusion from
two seeds; a systematic multi-seed sweep is left as future work.

\section{Ablations and analysis}
\label{sec:ablations}

\subsection{Inference-time knockout on trained checkpoints}
\label{sec:ablations:inference-ko}

The F4 knockout in \S\ref{sec:results:midscan} shows that zeroing the
eager term $a_t = M_{t-1} k_t$ during the \emph{probe's} state
reconstruction collapses the answer signal to chance. A stronger
inference-time test replaces the mixer's own forward-pass eager term
with zero on trained weights and evaluates the model on the actual
task. If the model has wound the eager term into its computation, the
prediction accuracy will not simply drop to chance ($0.25$); it will
be systematically \emph{wrong} because the weights are calibrated
around the perturbation.

\begin{table}[h]
\centering
\caption{Inference-time eager-KO on trained checkpoints, grouped by
seed category. Same model weights evaluated with the eager term
forced to zero in the mixer's forward pass. Chance $= 0.25$. Passing
\texttt{rev} seeds collapse by $\sim$30x to well below chance;
failing \texttt{rev} seeds show smaller but real degradation;
\texttt{fwd} seeds (passing and failing) collapse to $\lesssim
0.001$ hits under KO, exceeding the \texttt{rev} collapse. Group
means over $n$ seeds; per-seed data committed alongside.}
\label{tab:inference-ko}
\begin{tabular}{llccccc}
\toprule
Group & $n$ & Setting & $K_1{=}1$ & $K_1{=}2$ & $K_1{=}3$ & $K_1{=}4$ \\
\midrule
\texttt{rev} passing & 11 & eager on  & 0.996 & 0.969 & 0.967 & 0.952 \\
                     &    & eager off & 0.042 & 0.038 & 0.034 & 0.032 \\
\midrule
\texttt{rev} failing & 9  & eager on  & 0.957 & 0.464 & 0.309 & 0.260 \\
                     &    & eager off & 0.092 & 0.074 & 0.065 & 0.059 \\
\midrule
\texttt{fwd} passing & 1  & eager on  & 0.999 & 0.968 & 0.980 & 0.981 \\
                     &    & eager off & 0.000 & 0.000 & 0.000 & 0.000 \\
\midrule
\texttt{fwd} failing & 9  & eager on  & 0.950 & 0.478 & 0.322 & 0.269 \\
                     &    & eager off & 0.016 & 0.005 & 0.002 & 0.000 \\
\bottomrule
\end{tabular}
\end{table}

The KO accuracy for passing \texttt{rev} seeds is not just below chance
--- it is an order of magnitude below --- meaning the model does not
degrade to a random-token prior when the eager term is removed; it
emits confidently wrong tokens. The eager perturbation is the
computational pathway the model routes through, not a small correction
term. The \texttt{fwd} collapse to $\lesssim 0.001$ hits (both for the
one passing seed and for every failing seed) is stronger still: even
the small hop-1 competence the failing \texttt{fwd} baselines display
(\claim{$k_1 \sim 0.94$}) is downstream of the eager term, and without
it the model produces essentially no correct answer at any rung.

\subsection{Training-time causal control: eager term removed at training}
\label{sec:ablations:no-eager}

The inference-time knockout above shows that trained weights depend
critically on the eager term at inference. A separate question is
whether the mechanism can \emph{form at all} without the eager term.
We train a paired ablation under matched protocol (joint mode, 5000
steps, $n=5$ seeds; \texttt{rev} and \texttt{fwd} reduce to the same
model when the eager term is removed, and per-seed training was
byte-identical, see caption below) with the eager term disabled at
every step (\texttt{k\us{}eff} $= k$, standard fast-weight write with
no address perturbation).

\begin{table}[h]
\centering
\caption{Ablation: eager term removed at training time (joint curriculum,
$n=5$ seeds). Compare with Table~\ref{tab:ladder-joint} (same curriculum,
eager term on). Without the eager term the \texttt{rev} and \texttt{fwd}
variants collapse to the same model (both reduce to
\texttt{k\us{}eff} $= k$, a standard fast-weight write); we verified
this by running both --- final accuracies are byte-identical per seed.
$K_1{=}1$ single-hop retrieval works perfectly on every seed
(the outer-product write suffices), but chain composition to $K_1{=}4$
plateaus far below the eager-on ceiling: one of five seeds reaches
$\text{acc} = 0.503$ ($\sim$2$\times$ chance), the other four stall at
chance. The eager address-perturbation is what lifts chain composition
from $\sim$$0.50$ (partial residual-based pathway) to $\sim$$0.97$
(eager-mediated pathway).}
\label{tab:no-eager}
\begin{tabular}{lcccc}
\toprule
Seed & $K_1 = 1$ & $K_1 = 2$ & $K_1 = 3$ & $K_1 = 4$ \\
\midrule
0 & 1.000 & 0.660 & 0.566 & \textbf{0.503} \\
1 & 1.000 & 0.507 & 0.303 & 0.261 \\
2 & 1.000 & 0.472 & 0.266 & 0.272 \\
3 & 1.000 & 0.479 & 0.276 & 0.282 \\
4 & 1.000 & 0.422 & 0.273 & 0.255 \\
\bottomrule
\end{tabular}
\end{table}

\subsubsection*{Mechanism of the partial-pass ablated seed}

Probing the one partial-passing ablated seed (s0, $k_4{=}0.503$) with
F3 and F4 reveals a mechanism qualitatively different from the
eager-on passing seeds of \S\ref{sec:results:mechanism}:

\begin{itemize}
\item \textbf{F3 (logit lens).} L0=$0.250$, L1=$0.249$, L2=$0.470$,
      L3=$0.492$. Only the final two residuals lift above chance, and
      only to $\sim$$0.49$, matching the seed's $k_4$ accuracy.
\item \textbf{F4 (mid-scan buildup, ATOK answer signal).} L=0
      ATOK$={}0.268$; L=1 ATOK$={}0.379$; L=2 ATOK$={}0.211$. Weak
      signal at L=1, no clean per-layer buildup. The eager-mediated
      pathway that gives the sharp L=0-or-L=1 buildup in the
      eager-on passing seeds is absent.
\end{itemize}

The residual-based pathway that survives without the eager term is
weaker and more diffuse. It reaches $\sim$$0.5$ where the
eager-mediated pathway reaches $\sim$$0.97$, and it does not localise
to a single mixer's write. This is consistent with the interpretation
that the eager term is not a small booster but the primary
computational channel for single-pass chain composition on this task.

\subsection{Other planned ablations (future work)}

\begin{itemize}
\item vs standard DeltaNet ($k, v - M^\top k$) as a canonical
      delta-family baseline
\item vs Mamba2 as a non-fast-weight strong SSM baseline
\end{itemize}

\section{Related work}
\label{sec:related}

\textbf{The lineage.} The primitive we study sits inside the
fast-weight programmer family, which runs from
\citet{schmidhuber1992fastweights} through
\citet{schlag2021linear,schlag2021deltanet} to the recent DeltaNet
work of \citet{yang2024parallelizing,yang2024gated}. All of these
share the outer-product write $M \mathrel{{+}{=}} \beta_t \, k_t' \,
v_t'^{\top}$; they differ in how $k_t'$ and $v_t'$ depend on the
current state. Delta-family mixers modify the \emph{value} being
written (the residual $v_t - M^{\top} k_t$); FENRIR modifies the
\emph{key} of the write via $k_{\text{eff},t} = k_t + M_{t-1} k_t$
(reverse-lookup variant) or $k_t + M_{t-1}^{\top} k_t$
(forward-lookup variant).

\textbf{On novelty of the primitive.} We do not claim strong novelty
for the FENRIR update rule. Prior work on fast-weight programmers,
self-referential weight matrices \citep{schlag2021deltanet}, and
state-augmented linear-attention variants has explored a large space
of state-dependent write dynamics, and it is plausible that an update
functionally equivalent to FENRIR has appeared under another name. We
searched the literature we know
(\S\ref{sec:related:closest} below) and did not find an exact match;
we welcome corrections. What we claim in this paper is not the
primitive but a mechanistic account of how one instance of this
family learns (and often fails to learn) 2-hop chain composition
inside a single forward pass.

\subsection{Closest prior art we found}
\label{sec:related:closest}

The DeltaNet family
\citep{schlag2021deltanet,yang2024parallelizing,yang2024gated} is
structurally closest: same outer-product write, same forward read,
same $\mathcal{O}(d^2)$ state, similar chunk-parallel kernel for the
reverse-lookup variant. The intervention is in a different place
(value vs key). GatedDeltaNet \citep{yang2024gated} adds scalar
gating on the residual, which is compatible with FENRIR's address
perturbation in principle.

We searched several other lines of fast-weight and linear-attention
work for a state-dependent modification to the write key before the
outer product. The closest analogs we found each apply their
state-dependent term in a different place than FENRIR does:

\begin{itemize}
\item \citet{ba2016fastweights} iterate a fast-weight read
      $h^{s+1} = \mathrm{LN}(W h_{\text{prev}} + C x + A h^s)$
      inside the recurrent step: the fast-weight matrix $A$ is read
      with the current hidden state and the result is added back into
      that state. The state-dependent additive is structurally
      similar to FENRIR's but strictly on the READ side; $A$ is not
      built afterwards from the modified vector.
\item \citet{ccq2026} apply $q_{\text{clean}} =
      (I - \lambda \Sigma) q$ where $\Sigma$ is a running covariance
      of past keys, a state-dependent additive correction to the
      QUERY on the READ side. Mirror image of FENRIR's intervention,
      not the same one.
\item \citet{qdelta2026} build a mixed vector $k + \lambda q$ that
      appears in the erase term $(I - \beta (k k^{\top} + \lambda q
      k^{\top}))$; still not a state-dependent addition to $k$ before
      the write.
\item Self-referential weight matrices \citep{irie2022srwm} let a
      fast-weight matrix generate its own $q$, $k$, $v$, $\beta$
      projections from the current state, then apply a delta update.
      Self-reference is on GENERATION of the projections, not on a
      $k + M k$ residual after projection.
\end{itemize}

We surveyed 15+ additional variants (DeltaProduct
\citep{deltaproduct2025}, MesaNet, Longhorn, Titans, TTT-Linear
\citep{sun2024ttt}, Kaczmarz Linear Attention, and the wider
recurrent-linear-attention family) and did not find a closer match.

The two closest self-referential fast-weight families were verified
against FENRIR's update rule directly:

\begin{itemize}
\item Recurrent Fast Weight Programmers \citep{irie2021rfwp} include
      a Recurrent Delta Net (RDN) whose $k$, $v$, $q$, $\beta$
      projections are augmented with a state-dependent term
      $R_{\bullet} \tanh(y^{(i-1)})$ derived from the previous output.
      This is state-dependent only through the previous output vector
      $y^{(i-1)}$ (a $d$-dim quantity routed through additional slow
      weights); no term of the form $M k$ enters the write. Not the
      FENRIR update.
\item The Modern Self-Referential Weight Matrix \citep{irie2022srwm}
      generates its own projections $k_t = [W_{t-1}]_k \phi(x_t)$
      from the current fast-weight matrix (i.e.\ the state IS the
      projection matrix), then applies a standard delta write
      $W_t = W_{t-1} + \sigma(\beta_t)(v_t - \bar{v}_t) \otimes
      \phi(k_t)$. FENRIR uses a fixed slow-weight projection
      $k_t = W_k x_t$ and adds an explicit state-dependent
      perturbation $k_{\text{eff}} = k_t + M_{t-1} k_t$ before the
      write. Different level of self-reference; not the same update.
\end{itemize}

We were unable to find a match for FENRIR's specific update rule in
the fast-weight or linear-attention literature we surveyed. We do
not claim exhaustive coverage of obscure workshops, non-arXiv venues,
or in-progress preprints, and welcome citations we missed. The
paper's argument (formation-mode bimodality, mid-scan mechanism,
causal knockout) does not rest on primitive novelty.

\subsection{Chain composition and mechanism study}

\textbf{Composition benchmarks.} Grokked-Transformers
\citep{wang2024grokked} argues dense transformers reach reliable
multi-hop composition only after extended training and frames the
capability as a target for architectural inductive biases. Zoology
\citep{arora2023zoology} and BASED \citep{arora2024based} isolate
single-hop associative recall as a linear-attention weak spot; they
do not probe multi-hop composition inside a single forward pass. Our
task (injective 2-hop chains with truncated supervision) is
intentionally minimal: it forces the model to link two facts through
a shared bridge in one pass, with a clean $1/K_2$ chance floor.

\textbf{Mechanism probes.} Logit-lens
\citep{nostalgebraist2020logitlens} inspired our F3 layer
localisation probe. Induction-head circuit analysis
\citep{olsson2022incontext} inspired the paired-knockout pattern we
use in F4. Both techniques were originally applied to attention
residual streams; we adapt them to the fast-weight state directly.
Formation-mode bimodality is a phenomenon adjacent to grokking
\citep{wang2024grokked} in the sense that a fraction of training
runs pass a capability transition and the rest do not; we do not
attempt to bridge the two literatures here.

\textbf{State-space baselines.} Mamba and Mamba2
\citep{gu2023mamba,dao2024mamba2} do not carry a fast-weight state
and are not intended to compete on tasks that require
in-context storage-and-retrieval. We treat Mamba2 as a swap-in
architectural baseline in the planned ablations, not as a competitor.

\section{Limitations}
\label{sec:limitations}

\begin{itemize}
\item \textbf{Formation bimodality.} At $n = 5$ seeds under mixed
      training the \texttt{rev} variant passes $2/5$. We do not have a
      training-time diagnostic that predicts which seeds will pass, and
      we do not offer a recipe that raises the pass rate.
\item \textbf{Task scope.} Every claim in this paper is about the 2-hop
      injective truncated composition task at $K_1 \in \{1,\dots,4\}$,
      $K_2 = 4$, at $d_{\text{model}} = 256$, $d_{\text{key}} = 32$,
      $3$ FENRIR blocks, on synthetic data. Whether the mechanism
      survives naturalistic data or deeper chains is out of scope.
\item \textbf{Variant coverage.} We report two operator directions
      (\texttt{rev}, \texttt{fwd}) and two training curricula (mixed,
      joint). Other axes -- optimizer sensitivity, alternative
      scaffolds, learned $\alpha$-gating on the eager term -- are not
      ablated.
\item \textbf{No LM transfer here.} The paper argues about a
      composition mechanism on a controlled task; it does not
      demonstrate that FENRIR is a competitive language-model mixer at
      any scale.
\end{itemize}

\section{Conclusion}
\label{sec:conclusion}

Chain composition inside a single mixer forward pass is attainable but
formation-mode contingent: on the address-perturbed fast-weight mixer we
study, a fraction of seeds train to full $K_1{=}4$ accuracy and the rest
stall at chance. The passing seeds carry a mid-scan mechanism signature
that a knockout of the address-perturbation term causally destroys. The
reproduction code and audit workflow are packaged so future
extensions can be checked against the convention tests before any new
claim leaves the repository.

\bibliographystyle{plainnat}
\bibliography{references}

\appendix

\section{Chunk-parallel kernel derivation (rev variant)}
\label{app:chunked}

The \texttt{rev} recurrence at position $t$ is
\begin{align}
a_t &= M_{t-1} k_t, & w_t &= \beta_t (k_t + a_t), \\
M_t &= M_{t-1} + w_t v_t^{\top}, & y_t &= q_t^{\top} M_t .
\end{align}
Fix a chunk of size $c$ with entry state $M_0$. For a position $t$ inside
the chunk, $M_{t-1} = M_0 + \sum_{s<t} w_s v_s^{\top}$, so
\begin{equation}
w_t = \beta_t \Big( k_t + M_0 k_t + \sum_{s<t} (v_s^{\top} k_t)\, w_s \Big).
\end{equation}
Stack the rows $w_t^{\top}$ into $W \in \mathbb{R}^{c \times d}$ and the keys
$k_t^{\top}$ into $K \in \mathbb{R}^{c \times d}$. The coupled equations above
become a single linear system,
\begin{equation}
\big(I - \operatorname{diag}(\beta)\, T\big)\, W
  = \operatorname{diag}(\beta)\,(K + P),
\qquad
P := K M_0^{\top}, \quad
T_{ts} := v_s^{\top} k_t \ \ (s < t),
\end{equation}
where $T$ is strictly lower triangular, so $I - \operatorname{diag}(\beta)\, T$
is unit lower triangular. $W$ is therefore recovered by one batched triangular
solve per chunk (\texttt{torch.linalg.solve\_triangular}), with no
per-position Python loop. The chunk output and exit state then follow as
\begin{equation}
Y = Q M_0 + \operatorname{tril}\!\big(Q W^{\top}\big)\, V,
\qquad
M_c = M_0 + W^{\top} V,
\end{equation}
where the $\operatorname{tril}$ includes the diagonal (read-after-write). The
per-chunk cost is $O(d^2 c + d c^2 + c^3)$ and the serial depth is $L/c$
instead of $L$.

Only the \texttt{rev} variant admits this triangular closure. The
\texttt{fwd} eager term $M^{\top} k$ produces a bilinear reduction over the
past, $\sum_{s<t} (w_s^{\top} k_t)\, v_s$, that does not collapse to a single
triangular solve without further algebra, so the \texttt{fwd} variant falls
back to the sequential path in \texttt{model.py}; a chunk-parallel
\texttt{fwd} form is left as an engineering follow-up. The state $M$ is
accumulated in \texttt{fp32} regardless of the caller's dtype, because bf16
accumulation of this long-horizon sum degrades retrieval; a
chunked-versus-sequential parity test (\texttt{tests/test\_parity.py}) pins
the two paths to the same output.

\section{Audit workflow and the operator-mismatch bug class}
\label{app:audit}

The mechanism probes (Section~\ref{sec:setup:probes}) reconstruct the mixer
state $M_t$ outside the module so that individual terms can be knocked out.
That reconstruction is only trustworthy if it uses the \emph{same} operator
the mixer used. The failure mode we guard against is an external einsum whose
contraction axis does not match the mixer's storage convention: under the
convention $M[i,j] = k_{\text{eff}}[i]\, v[j]$, the two eager terms
$M k$ and $M^{\top} k$ differ only by which axis is contracted, so a probe
that picks the wrong axis runs without error and returns plausible numbers
while silently measuring a different operator than the model was trained
with. On this project that class of bug was caught twice before it reached a
reported number.

The safeguard we adopted has two parts. First, the mixer exposes its
per-step eager term under a variant-agnostic attribute, so a probe reads the
mixer's own operator rather than recomputing one and risking a mismatch.
Second, every storage-convention and per-variant operator claim is pinned as
a \texttt{pytest} assertion in \texttt{code/tests/test\_conventions.py}, so a
convention change breaks CI rather than a result. The passing test suite is
the audit. The internal audit document that originally motivated this
discipline is publisher-internal and is not shipped with the reproduction
repository.

\end{document}
